% =====================================================================
% RAG Architecture Portfolio � Completion Status Report
% Compile: pdflatex portfolio_status.tex  (twice for references)
% =====================================================================
\documentclass[11pt,a4paper]{article}

\usepackage[T1]{fontenc}
\usepackage[utf8]{inputenc}
\usepackage{lmodern}
\usepackage[a4paper,margin=2.1cm]{geometry}
\usepackage{parskip}
\usepackage{microtype}

\usepackage{amsmath,amssymb}
\usepackage{booktabs}
\usepackage{longtable}
\usepackage{tabularx}
\usepackage{array}
\usepackage{enumitem}
\setlist[itemize]{leftmargin=1.3em,itemsep=1.5pt,topsep=1.5pt}
\setlist[enumerate]{leftmargin=1.5em,itemsep=1.5pt,topsep=1.5pt}

\usepackage{xcolor}
\usepackage[most]{tcolorbox}
\newtcolorbox{notebox}{
  colback=gray!6, colframe=gray!55!black,
  boxrule=0.4pt, arc=2pt, left=6pt, right=6pt, top=4pt, bottom=4pt
}

\usepackage{ragged2e}
\usepackage{fancyhdr}
\pagestyle{fancy}
\fancyhf{}
\fancyhead[L]{\small RAG Architecture Portfolio --- Completion Status}
\fancyhead[R]{\small\thepage}
\renewcommand{\headrulewidth}{0.3pt}

\usepackage[hidelinks]{hyperref}
\hypersetup{
  pdftitle={RAG Architecture Portfolio Completion Status},
  pdfauthor={Project Owner},
  bookmarksnumbered=true
}

\usepackage{titlesec}
\titleformat{\section}{\large\bfseries}{\thesection}{0.6em}{}
\titleformat{\subsection}{\normalsize\bfseries}{\thesubsection}{0.5em}{}
\titlespacing*{\section}{0pt}{1.0em}{0.35em}
\titlespacing*{\subsection}{0pt}{0.7em}{0.25em}

\usepackage{makecell}
\renewcommand\theadfont{\bfseries\small}
\renewcommand\cellalign{tl}

\newcommand{\code}[1]{\texttt{#1}}

% Table macros for tight cells
\newcolumntype{L}[1]{>{\RaggedRight\arraybackslash}p{#1}}
\newcolumntype{C}[1]{>{\Centering\arraybackslash}p{#1}}

% =====================================================================
\begin{document}

\begin{center}
{\LARGE\bfseries RAG Architecture Portfolio}\\[4pt]
{\large Completion Status Report}\\[6pt]
\rule{\linewidth}{0.4pt}
\end{center}

\vspace{-0.4em}

\begin{notebox}
\textbf{Scope.} This report records the terminal status of the eleven
retrieval architectures audited under the present program. It
describes, for each architecture, what was feasible, what was
completed, and where the work stopped. No empirical retrieval result
is reported here. All statuses reflect work completed at the time of
writing.
\end{notebox}

% =====================================================================
\section{Overview}

Eleven retrieval architectures were considered. Each was audited for
retrieval-only execution under three binding constraints: academic
correctness, comparative fairness, and architectural integrity.
Two hard resource constraints also applied: no hosted LLM API and no
persistent GPU.

Of the eleven architectures:

\begin{itemize}
  \item \textbf{Two} reached full completion and are ready for
    dataset-level execution.
  \item \textbf{One} reached partial execution and is retained with a
    declared structural limitation.
  \item \textbf{Eight} stopped at a well-defined blocker and are
    recorded as deferred or excluded with stated reasons.
\end{itemize}

The following two sections summarise each architecture and then
tabulate the terminal status, the reason for stopping, and whether
the vendor's original architecture or the paper's architecture was
altered in any way.

% =====================================================================
\section{Per-Architecture Status}

\subsection*{\code{my\_rag} (Project-Owned Reference)}

Fully implemented and validated in the infrastructure pilot. The
pipeline covers canonical input, structural parsing, hierarchical
child chunking, dense and sparse embedding with BGE-M3, Qdrant
persistence, hybrid RRF retrieval, parent aggregation, and common
evaluation. No LLM is used at any stage. \textbf{Status: complete.
Not yet executed on any real dataset.}

\subsection*{\code{longrag} (Controlled Adaptation)}

Long-unit construction, dense embedding, exact reference index,
Qdrant storage, retrieval, canonical expansion, and evaluation-case
construction all passed synthetic acceptance. Storage equivalence
against an exact reference index was validated with top-$k$ overlap
of $1.0$ at all cutoffs. Embeddings have been generated in a Colab
T4 session. \textbf{Status: complete.}

\subsection*{\code{pathrag} (Paper-Derived Comparator)}

The foreign repository is unlicensed, so only a paper-derived
clean-room comparator is executable. The comparator implements graph
construction, bounded simple-path enumeration, resource propagation,
path scoring, path ranking, neutral textualisation, and canonical
document expansion. The comparator operates without embedding
augmentation. An attempt to augment the comparator with BGE-M3
embeddings failed due to a combination of licensing and interface
constraints; this is documented in a separate failure report.
\textbf{Status: partial. Retained as a structural comparator with a
declared limitation.}

\subsection*{\code{raptor} (Controlled Adaptation)}

Leaf chunking, tree serialisation, and ranked-node retrieval were
implemented. The architecture-defining summarisation stage requires
an LLM at index time. A frozen Qwen2.5-7B-Instruct request bundle was
prepared for a Colab T4 pilot. The full summarisation stage was not
executed because the number of LLM calls required for a long-context
dataset is not sustainable without hosted infrastructure.
\textbf{Status: stopped at the summarisation stage.}

\subsection*{\code{hipporag2} (Controlled Adaptation)}

Graph construction, personalised PageRank, passage ranking, and
canonical mapping were implemented. The architecture-defining OpenIE
stage requires an LLM at index time. Graph construction for a single
long-context dataset requires on the order of $10^{7}$ LLM tokens.
This is not sustainable without hosted infrastructure.
\textbf{Status: stopped at the OpenIE stage.}

\subsection*{\code{lightrag} (Controlled Adaptation)}

Chunking, graph construction, and dual-level retrieval were
implemented. The architecture-defining entity and relation extraction
stage requires an LLM at index time. Of the three graph-based systems
considered, LightRAG has the highest graph-construction cost,
exceeding $10^{7}$ LLM tokens per dataset. This is not sustainable
without hosted infrastructure. \textbf{Status: stopped at the
extraction stage.}

\subsection*{\code{disretrieval} (Deferred)}

The paper and repository were audited statically. The repository
carries no software license. The architecture-defining neural RST
parser and its bundled checkpoint cannot be executed, imported,
modified, or redistributed under default copyright. Substituting the
parser with a heuristic or a generic constituency parser would
materially alter the architecture. The architecture is therefore
closed as deferred with reason.
\textbf{Status: deferred. Blocked on license.}

\subsection*{\code{memorag} (Excluded)}

The repository is Apache-2.0 licensed, but the architecture requires
a GPU with 16--24~GB VRAM at both index time and query time, and
requires an LLM at both stages. Neither requirement is available on
the current host. \textbf{Status: excluded on hardware grounds.}

\subsection*{\code{kag} (Excluded)}

The repository is Apache-2.0 licensed, but the architecture requires
a multi-service environment (OpenSPG, Neo4j) together with hosted LLM
access. The infrastructure is not reproducible within the scope of
the study. \textbf{Status: excluded on infrastructure grounds.}

\subsection*{\code{metadata\_rag} (Paper-Derived, Planned)}

No foreign repository exists. The paper is a preprint under revision.
A project-owned clean-room implementation is feasible and would
require approximately three weeks of part-time work. Query metadata
inference can be implemented with rules rather than an LLM.
\textbf{Status: planned, not started.}

\subsection*{\code{hikey} (Excluded)}

The architecture requires multimodal inputs (PDF, tables, images,
layout) and a GPU. No public repository is available. The canonical
envelope used across the program is text-only.
\textbf{Status: excluded on input-contract grounds.}

% =====================================================================
\section{Terminal Status Table}

Table~\ref{tab:status} records the terminal status of each
architecture, the reason for stopping where applicable, and whether
the vendor's original code or the paper's architecture was altered.

\begin{center}
\small
\begin{longtable}{@{}L{2.5cm} L{1.9cm} L{3.5cm} L{3.1cm} L{3.1cm}@{}}
\caption{Terminal status of the eleven audited architectures.}
\label{tab:status}\\
\toprule
\thead{Architecture} & \thead{Status} & \thead{Reason for stop} &
\thead{Vendor architecture alteration} &
\thead{Paper architecture alteration} \\
\midrule
\endfirsthead
\toprule
\thead{Architecture} & \thead{Status} & \thead{Reason for stop} &
\thead{Vendor architecture alteration} &
\thead{Paper architecture alteration} \\
\midrule
\endhead
\bottomrule
\endfoot

\code{my\_rag} &
Complete &
--- &
None (project-owned). &
None. \\
\addlinespace

\code{longrag} &
Complete &
--- &
Controlled adaptation. Foreign code not executed; clean-room adapter only. &
Architecture-defining logic preserved. Embedder replaced with BGE-M3 (labelled). \\
\addlinespace

\code{pathrag} &
Partial &
Embedding augmentation failed. Comparator kept without embeddings. &
Foreign code not executed (unlicensed). Clean-room comparator only. &
Algorithmic logic preserved. Embedding stage absent by design. \\
\addlinespace

\code{raptor} &
Stopped &
Index-time LLM summarisation cost not sustainable. &
Not executed. Adapter contracts only. &
Architecture preserved in contracts. Summarisation stage not run. \\
\addlinespace

\code{hipporag2} &
Stopped &
Index-time LLM OpenIE cost not sustainable. &
Not executed. Adapter contracts only. &
Architecture preserved in contracts. OpenIE stage not run. \\
\addlinespace

\code{lightrag} &
Stopped &
Index-time LLM extraction cost not sustainable. &
Not executed. Adapter contracts only. &
Architecture preserved in contracts. Extraction stage not run. \\
\addlinespace

\code{disretrieval} &
Deferred &
Architecture-defining RST parser has no license. &
Not executed. No code touched. &
Not altered. Deferred rather than substituted. \\
\addlinespace

\code{memorag} &
Excluded &
Requires 16--24~GB VRAM GPU and LLM at index and query time. &
Not executed. &
Not altered. Out of scope. \\
\addlinespace

\code{kag} &
Excluded &
Requires OpenSPG and Neo4j service environment and hosted LLM. &
Not executed. &
Not altered. Out of scope. \\
\addlinespace

\code{metadata\_rag} &
Planned &
No repository. Clean-room implementation not started. &
Not applicable (no vendor code). &
Not altered. Implementation not started. \\
\addlinespace

\code{hikey} &
Excluded &
Requires multimodal input contract and GPU. No public repository. &
Not executed. &
Not altered. Out of scope. \\
\bottomrule
\end{longtable}
\end{center}

% =====================================================================
\section{Preservation Summary Table}

Table~\ref{tab:preservation} summarises, in a compact form, which
architectures were executed, which were adapted, and which were
preserved without alteration. This view supports the fairness claim
of the comparison.

\begin{center}
\small
\begin{tabularx}{\linewidth}{@{}L{3.0cm} C{1.6cm} C{1.6cm} C{1.6cm} L{4.5cm}@{}}
\caption{Execution and preservation status.}
\label{tab:preservation}\\
\toprule
\thead{Architecture} &
\thead{Executed} &
\thead{Adapted} &
\thead{Vendor\ arch.\ intact} &
\thead{Note} \\
\midrule
\endfirsthead
\toprule
\thead{Architecture} &
\thead{Executed} &
\thead{Adapted} &
\thead{Vendor\ arch.\ intact} &
\thead{Note} \\
\midrule
\endhead
\bottomrule
\endfoot

\code{my\_rag}       & Yes & No  & N/A & Project-owned reference. \\
\code{longrag}       & Yes & Yes & Yes & Controlled adaptation, embedder labelled. \\
\code{pathrag}       & Yes & Yes & Yes & Paper-derived comparator, no embeddings. \\
\code{raptor}        & No  & No  & Yes & Contracts frozen; summarisation not run. \\
\code{hipporag2}     & No  & No  & Yes & Contracts frozen; OpenIE not run. \\
\code{lightrag}      & No  & No  & Yes & Contracts frozen; extraction not run. \\
\code{disretrieval}  & No  & No  & Yes & Deferred on license. \\
\code{memorag}       & No  & No  & Yes & Excluded on hardware. \\
\code{kag}           & No  & No  & Yes & Excluded on infrastructure. \\
\code{metadata\_rag} & No  & No  & N/A & Planned. \\
\code{hikey}         & No  & No  & Yes & Excluded on input contract. \\
\bottomrule
\end{tabularx}
\end{center}

% =====================================================================
\section{Notes on Interpretation}

\begin{itemize}
  \item \textbf{Preservation.} No architecture was executed with an
    unlabelled substitution of an architecture-defining component.
    Where a controlled substitution was performed, it is recorded as
    such. Where no substitution was possible, the architecture was
    stopped or deferred rather than approximated.
  \item \textbf{Vendor code.} No foreign repository was executed,
    imported, or modified. All executable adapters are project-owned.
  \item \textbf{Paper architecture.} The architecture-defining logic
    described in each paper is preserved in the corresponding
    contracts. Where the logic was not executed, the reason is
    recorded.
  \item \textbf{LLM cost.} The three LLM-heavy graph and tree
    architectures (RAPTOR, HippoRAG2, LightRAG) require hosted LLM
    calls at index time. Their aggregate cost, estimated in the range
    of $10^{7}$ tokens per dataset, is not sustainable within the
    study's resource envelope.
  \item \textbf{Fairness.} The remaining comparison set consists of
    architectures that can be executed under identical conditions:
    same canonical envelope, same embedder, same queries, same
    evaluation contract.
\end{itemize}

% =====================================================================
\section{Conclusion}

Two architectures reached full completion. One architecture is
retained with a declared structural limitation. Eight architectures
stopped at a well-defined blocker, each with an explicit reason. No
architecture-defining component was substituted without a label, and
no foreign code was executed, imported, or modified. The comparison
set is therefore internally clean, reproducible, and defensible under
academic review.

\end{document}